\documentclass[10pt]{amsart}
\usepackage[margin=1.8cm,top=1.7cm,bottom=1.9cm]{geometry}
\usepackage{amsmath,amssymb,booktabs}
\let\olditemize\itemize\renewcommand{\itemize}{\olditemize\setlength{\itemsep}{1pt}\setlength{\topsep}{2pt}\setlength{\parsep}{0pt}}
\setlength{\parskip}{2pt}
\newtheorem{theorem}{Theorem}\newtheorem{lemma}[theorem]{Lemma}\newtheorem{proposition}[theorem]{Proposition}
\newtheorem{corollary}[theorem]{Corollary}
\theoremstyle{definition}\newtheorem{definition}[theorem]{Definition}\newtheorem{remark}[theorem]{Remark}
\newcommand{\Z}{\mathbb Z}\newcommand{\Q}{\mathbb Q}\newcommand{\F}{\mathbb F}
\newcommand{\CP}{\mathbb{CP}}\newcommand{\CPb}{\overline{\mathbb{CP}}{}^2}
\newcommand{\PD}{\operatorname{PD}}\newcommand{\CS}{\operatorname{CS}}
\newcommand{\V}{\textsf{[V]}}\newcommand{\Pl}{\textsf{[P]}}\newcommand{\Sp}{\textsf{[S]}}
\makeatletter\def\section{\@startsection{section}{1}{\z@}{.7\baselineskip}{.3\baselineskip}{\normalfont\bfseries}}\makeatother
\title{A classical reading of the cosmetic-surgery argument}
\date{}
\begin{document}
\maketitle
\vspace{-2em}
\noindent\textit{Status.} This note reorganizes the manuscript's argument in the language of Floer, Donaldson, Kronheimer--Mrowka, Feehan--Leness and DLME. It has not been checked against the manuscript line by line. Labels: \V\ means checked by hand or by machine; \Pl\ means plausible; \Sp\ means speculative. Supporting notes are in \texttt{notes/}. Throughout, $Y_r=S^3_r(K)$, and DLME means Daemi--Lidman--Miller Eismeier, arXiv:2410.21248.

\section{The theorem and the shape of the argument}
By Hanselman and DLME (Cor.~1.4), a purely cosmetic pair must be $\{\pm2\}$ with $g(K)=2$. The method proves the following more symmetric statement (\texttt{notes/T4}).

\begin{theorem}[conditional on Theorem~\ref{thm:FL}]\label{thm:main}
Let $K\subset S^3$ be nontrivial. There do not exist negative-definite cobordisms $V\colon Y_{-2}\to Y_2$ and $V'\colon Y_2\to Y_{-2}$ with $b_1=0$, both inducing isomorphisms on $I(\,\cdot\,;\F_2)$, such that $\pi_1(V')$ is normally generated by the image of one end. In particular the purely cosmetic surgery conjecture holds.
\end{theorem}

\noindent The argument has three steps.
\begin{enumerate}
\item \emph{Algebra.} Surgery exact triangles yield integral maps $B\colon I(Y_2)\to I(Y_{-2})$ and $H\colon I(Y_{-2})\to I(Y_2)$ that are isomorphisms mod 2. Together with $V$ they give an automorphism $w$ of the $2$-adic lattice of $I(Y_2)$. Relative invariants of two caps give a nonzero pairing.
\item \emph{Nonvanishing.} A suitable power $w^M$ is the identity mod $2^N$. So a \emph{secondary} (family-of-metrics) relative Donaldson invariant $\Omega(X_M)$ of the capped composite is $\equiv q\not\equiv0\pmod{2^N}$.
\item \emph{Vanishing.} Once the coupled Dirac index satisfies $n_D\ge1$, an $\mathrm{SO}(3)$-monopole (Pidstrigach--Tyurin, Feehan--Leness) relation expresses $2^{n_D-1}\Omega$ through Seiberg--Witten-type reducibles. A family adjunction inequality shows there are none. Hence $\Omega=0$, a contradiction.
\end{enumerate}
The cosmetic map is used to make the (negative-definite) cobordism $V$ compose with itself. This gives many $(-4)$-spheres per unit of $b^+$.

\section{Surgery triangles}\label{sec:triangles}
We use Floer's irreducible instanton homology with determinant-one gauge group. $Y_{\pm1}$, $Y_{\pm2}$ and $S^3$ carry only central reducibles, $I(S^3)=0$, and $Y_0$ carries $w$ with $w\cdot A$ odd. The following dictionary is in \texttt{notes/T5}; it is \V\ for the topology and \Pl\ for the analysis.
\begin{itemize}
\item $f_+\colon I^w(Y_0)\to I(Y_2)$ and $f_-\colon I(Y_{-2})\to I^w(Y_0)$ are composites of Floer-triangle maps. They are isomorphisms because the $S^3$ term vanishes.
\item $g_\pm$ is DLME's family map $g_1$ for the triads $\{2,\infty,0\}$ and $\{0,\infty,-2\}$, with middle term $I(S^3)\oplus I(S^3)=0$. Their $h$-maps give $f_+g_+\simeq1$ and $g_-f_-\simeq1$.
\item \emph{The distance-four map.} Let $W'=(-X_2(K))\cup_{S^3}X_{-2}(K)\colon Y_2\to Y_{-2}$. Then $b^+=0$ and $H_2=\langle F_l,F_r\rangle\cong\langle-2\rangle^2$. The union of the two cores is a sphere $S=F_l-F_r$ with $S^2=-4$, and $\partial\nu S=L(4,1)$. Let $\Phi_0$ be the map defined by the one-parameter family of metrics running from the $S^3$-splitting to the $L(4,1)$-splitting, with Donaldson's class $\mu(S)$ inserted.
\end{itemize}

\begin{proposition}[\texttt{notes/T2}, \texttt{notes/T5}]\label{prop:B}
\
\begin{enumerate}
\item The class $\PD(S)=2y$ is even. So the bundles $c$ and $c+\PD(S)$ are the same $\mathrm{SO}(3)$-bundle, with both end identifications twisted by the flat $\Z/2$ character. Writing $\iota$ for this twist, the map
\[
B=\Phi_0-s\,\iota\Phi_0\iota
\]
is an integral chain map. The $L(4,1)$-wall (minimal trace-zero cap, $\kappa=\tfrac14$) lies in the opposite eigenspace of $\iota$-conjugation and cancels \Pl.
\item The insertion of $\mu(S)$ is forced. The minimal $L(4,1)$-cap has framed index $2$, against $0$ for the $\mathbb{RP}^3$-cap of DLME; equivalently, it is forced by degree matching with $g_-g_+$ \V.
\item $B\simeq g_-g_+=(f_+f_-)^{-1}=:H^{-1}$ on homology mod 2, up to a degree-four twist \Pl, about 80\%.
\end{enumerate}
\end{proposition}

\begin{remark}
The manuscript's insertion (normalized holonomy over a sweep of loops across $S$ equal to $-1$) is $\pm\mu(S)$. The class is $c_1$ of the determinant line of $\bar\partial$ on $S\cong\mathbb P^1$, and its canonical section vanishes on the jumping-line divisor $V_S=\{E|_S\not\cong\mathcal O\oplus\mathcal O\}$ \V. The manuscript's zero-winding modification is an artifact of holonomy representatives: $V_S$ avoids every reducible with $v\cdot S=0$ automatically.
\end{remark}

\noindent\emph{Relation to DLME's remark on $\pm2$.} In DLME's bookkeeping, $B$ has Chern--Simons level minus degree$/8$ equal to $+\tfrac18-\eta$, which is unfavourable; DLME's $g_1$ at $\pm1$ has $-\tfrac18-\eta$. The difference is exactly the forced $\mu(S)$. So $B$ is DLME's ``trivial middle term, chain map, unfavourable filtration'' variation, and the Chern--Simons filtration alone cannot conclude \V.

\section{Relative invariants}
Take a symplectic closure $X_0\supset Y_0$ (Gabai, Eliashberg--Thurston, Eliashberg caps, Kronheimer--Mrowka) with nonzero Donaldson invariant $D^w_{X_0}(x^mh^a)$, $h\cdot A=1$. Insert copies of the bridge $T=f_-\circ V'\circ f_+$, which is a rational isomorphism of $I^w(Y_0)$. Carry the surface classes across $T$ by a nilpotent deformation over $\Q[t_i]/(t_i^2)$. By Cayley--Hamilton, some $T^k$ with $k\ge1$ preserves nonvanishing. Cutting along an interior $Y_2$ gives relative invariants
\[
\Psi_-\in I(Y_2;\Z),\qquad \Psi_+\in I(-Y_2;\Z),\qquad \langle\Psi_+,\Psi_-\rangle=q\neq0,
\]
from manifolds $C_\pm$ with $b^+>0$, each blown up once more \Pl. Genus two of $K$ is not used here; the tight genus-two surfaces come from the Lefschetz pencil \Pl.

\section{The secondary invariant and its nonvanishing}
Let $w=(V B)^3\circ V(HB)$, a self-cobordism of $Y_2$ built from four copies of $W'$, three copies of $V$ and one copy of $W_H$. Let $X_M$ be the closed manifold obtained by capping $C_-\cup(\text{pads})\cup w^M\cup(\text{pads})\cup C_+$.
\begin{itemize}
\item \emph{Topology} \V. $X_M$ contains $n\approx4M$ disjoint $(-4)$-spheres $S_i$ and $m=M$ positive pieces. Each positive piece is $P\cong\CP^2\#5\CPb$ for every $K$, contains a square-zero sphere $U$, and contains the plumbing $(-S_j,U,S_{j+1})$ with squares $(-4,0,-4)$. Since all the $W'$ pass through $S^3$, $X_M$ is a connected sum along those spheres. Its ordinary Donaldson and Seiberg--Witten invariants therefore vanish.
\item \emph{The secondary invariant} (\texttt{notes/T1}, \Pl). $\Omega(X_M)$ is the count of instantons over the cube $[0,1]^n$ of metrics whose faces stretch the necks $S^3$ and $L(4,1)$, cut down by $\mu(S_1)\cdots\mu(S_n)$ and the cap classes, and summed with signs over the bundles $w_0+\sum e_i\PD(S_i)$. Single-bundle counts are not invariant; the signed sum is. By gluing, $\Omega(X_M)=\pm\langle\Psi_+,w^M\Psi_-\rangle$.
\item \emph{Nonvanishing} \V\ as algebra. Each factor of $w$ is an isomorphism mod $2$, so $w$ is invertible on $L\otimes\Z_{(2)}$, where $L=I(Y_2;\Z)/\mathrm{Tor}$. Choose $N$ with $2^N\nmid q$ and take $M$ divisible by the exponent of $\mathrm{GL}_r(\Z/2^N)$. Then $\Omega(X_M)\equiv\pm q\not\equiv0\pmod{2^N}$.
\end{itemize}

\section{Coupling to a spin$^c$ structure}
Fix a characteristic class $l$ and put $\Lambda=l+c$, $\Theta=(\Lambda^2-\sigma)/4$ and $\kappa=c_2-c^2/4$. The coupled Dirac operator on $W^+\otimes E$ has complex index $n_D=\Theta-\kappa$ on closed manifolds. For an instanton $A$ from $\alpha$ to $\alpha'$ on a cobordism, APS gives \V\ (\texttt{notes/B1})
\[
n_D(A)=\Theta_W-E(A)+\rho_t(\alpha)-\rho_{t'}(\alpha'),\qquad \rho_t(\alpha)=\tfrac12\big(\eta(D_{\alpha,t})-h\big)+\tfrac14\eta_{\rm sign}(Y).
\]
Hence $\CS-\rho_t$ is integer-valued up to a constant, and $8(\CS-\rho_t)-\mathrm{gr}$ is an integer grading from which $\CS$ cancels. The ``spin-filtered'' Floer group is really a \emph{spin-graded} one \V.

On $X_M$ the telescope gives $\Theta=m+\Theta_0$ and $8\kappa=n+3m+c$, so
\[
n_D(X_M)=\frac{5m-n}{8}+c_1 .
\]
Each copy of $W'$ lowers $n_D$ by $\tfrac18$ (one parameter, one degree-two insertion, $\Theta=0$). Each positive piece raises it by $\tfrac58$ ($b^+=1$ costs $\tfrac38$, and $\Theta=1$ gains $1$) \V.

\section{The secondary Feehan--Leness relation}
\begin{theorem}[\texttt{notes/T1}, Thm.~A; \Pl, about 70\%]\label{thm:FL}
Suppose $n_D\ge1$. After $n_D-1$ cuts by sections of the phase-weight-two line, the free part of the cut-down $\mathrm{SO}(3)$-monopole moduli space over the cube is a compact oriented $1$-manifold, possibly with rational multisection weights. Its ends are:
\begin{enumerate}
\item the instanton links $\CP^{n_D-1}$, each cut to $2^{n_D-1}$ points, with total $2^{n_D-1}\Omega$;
\item lens-face ends, which cancel in pairs of bundles $c$, $c+\PD(S_i)$;
\item ends at abelian limits: Seiberg--Witten classes $K=\Lambda+v$ at Feehan--Leness level $\ell(v)=\kappa+v^2/4\ge0$, broken limits, and mixed limits containing zero-spinor abelian segments.
\end{enumerate}
Ends at the $S^3$-faces without an abelian piece vanish for dimensional reasons. Hence $2^{n_D-1}\Omega=-\#\{\text{abelian ends}\}$.
\end{theorem}
\noindent\emph{Inputs.}
\begin{itemize}
\item Standard: Feehan--Leness compactness, transversality and instanton links; the Pidstrigach--Tyurin structure; Morgan--Mrowka--Ruberman and Donaldson for psc necks with stabilizers; weighted Stokes.
\item New: the lens cancellation carried into the coupled problem (the Dirac determinant is complex and adds no sign; about 70\%). The treatment of mixed limits by projecting away abelian zero-spinor segments (the manuscript's ``necessary projection''; 55--60\%) is the single genuinely new analytic idea. It replaces Feehan--Leness's obstructed gluing at reducibles.
\end{itemize}

\section{The Seiberg--Witten side is empty}
Three classical facts enter.
\begin{enumerate}
\item \emph{$\mu$-classes on reducibles} (\texttt{notes/T2}, \V). At $E=L_1\oplus L_2$ with $v=c_1(L_1)-c_1(L_2)$, $\mu(S)$ restricts to the link as $\pm\tfrac{v\cdot S}{2}\,h$, and $V_S$ contains the reducible if and only if $v\cdot S\neq0$. So an insertion is met only if $v\cdot S\ne0$, which is even. Otherwise it requires a bubble on $S$, i.e.\ a unit of Feehan--Leness level, with distinct bubbles for disjoint spheres.
\item \emph{$b^+=1$ chambers on rational-surface pieces} (\texttt{notes/T3}). Type-S solutions on $P$ solve the Seiberg--Witten equations for $K$ perturbed by $\beta^+\in 2\pi\Lambda$.
\begin{itemize}
\item A family of toric K\"ahler metrics with $\mathrm{Scal}\ge0$ (Guillemin, Abreu) with period $H=U+aX_1+bX_2$, $a,b\in[0,\tfrac14]$, gives by Weitzenb\"ock the band $(v\cdot H)(K\cdot H)<0$. Since $2\chi+3\sigma=4>0$ there are no walls at zero perturbation, and Kronheimer--Mrowka/Li--Liu wall-crossing shows the band is optimal \V.
\item Convexity over the vertex classes $H_I=U+\tfrac14\sum_{i\in I}X_i$ yields the chamber condition \V. Each $H_I$ is the period after rationally blowing down the $X_i$ for $i\in I$: the lens faces $a=\tfrac14$ are Wahl degenerations, and at a corner the piece is $\mathbb P(1,1,4)\to\CP^2$ \V.
\item Stretching along $S^2\times S^1$ gives $v\cdot U=0$, by the Kronheimer--Mrowka/Friedman--Morgan argument.
\item $\mathrm{Scal}\ge0$ is genuinely needed, because the argument requires empty moduli spaces, not vanishing counts.
\end{itemize}
\item \emph{Negative definiteness.} On negative pieces $v^2\le-\tfrac12(h^2+h'^2)$; with a general $V$ in place of $\phi$ this follows by Cauchy--Schwarz, and $\Lambda_V^2=\sigma_V$ is attainable because the closed piece is diagonal by Donaldson's theorem \V.
\end{enumerate}

\begin{lemma}[family adjunction inequality; \texttt{notes/T1} Lemma~3.1; \V\ as a lattice statement, and sharp]\label{lem:adj}
Let $\Gamma$ be a segment of consecutive pieces carrying an abelian configuration, including Uhlenbeck limits, separated lens caps and halves at cut necks. Let $s^-(\Gamma)$ be the number of insertions on $\Gamma$ not adjacent to positive pieces, and $m(\Gamma)$ the number of positive pieces. Then
\[
\kappa_\Gamma\ \ge\ \tfrac12\big(s^-(\Gamma)+m(\Gamma)\big).
\]
Globally, a class $v$ with $T(v)$ insertions met only through bubbles satisfies $-v^2/4+T(v)\ge\tfrac{n-m}2-C$.
\end{lemma}
\noindent Since $\kappa=\tfrac{n+3m}8+C'$, no Seiberg--Witten class has nonnegative level once $3n-7m>C_0$. Broken limits can isolate any segment, so a \emph{local} density condition is needed; one endpoint count gives the following.
\begin{itemize}
\item Broken fixed limits: excess $\ge(3k-7)m-3k+5$, so spacing $k\ge3$ suffices.
\item Mixed limits, priced by the $\mathrm{SO}(3)$-monopole index: excess $\ge(k-4)m-k+2$, so spacing $k\ge4$ is needed, and the bound is attained at $k=3$ \V.
\end{itemize}
This one lemma replaces the manuscript's fixed-run, mixed-run and component-score lemmas.

\section{Proof of Theorem~\ref{thm:main}}
Suppose $V$, $V'$ exist (for the cosmetic case, $V=\phi$ and $V'=\phi^{-1}$). Form $X_M$ with positive pieces spaced four apart and pads of length $p$, all chosen before $M$. Then
\[
n_D=\frac{M-p}{8}+c_1,\qquad 3n-7m=5M+3p .
\]
For $M$ large and divisible by the exponent of $\mathrm{GL}_r(\Z/2^N)$, we get $n_D\ge1$ and $3n-7m>C_0$. By Theorem~\ref{thm:FL} and Lemma~\ref{lem:adj}, $2^{n_D-1}\Omega(X_M)=0$, whereas $\Omega(X_M)\equiv\pm q\not\equiv0\pmod{2^N}$. This is a contradiction. The coefficients must be $\Z$ or $\Z_{(2)}$: over $\F_2$ the relation is empty once $n_D\ge2$, which is why $B$ has to be integral.

\emph{Where the cosmetic map enters.} We need $\tfrac15<m/n$ (so that $n_D\to\infty$) and $m/n<\tfrac37$ (so that the Seiberg--Witten side is empty), together with spacing $4$. Without $V$, every $B$ must be followed by $H$, because the reversed interval $Y_{-2}\to Y_2$ has $b^+=2$. That forces spacing $1$, and then Seiberg--Witten classes do contribute. This is consistent with $HB\simeq1$ mod 2.

\section{Assessment}
\begin{itemize}
\item \emph{Is it alien?} No. The ingredients are DLME and CDX triangles, Donaldson's $\mu$-classes, Kronheimer--Mrowka nonvanishing, Feehan--Leness/Pidstrigach--Tyurin localization, and $b^+=1$ wall-crossing on rational surfaces. The genuinely new idea is to use the $\mathrm{SO}(3)$-monopole cobordism as a \emph{vanishing theorem for a secondary invariant}, with the cosmetic map arranging the ratio of $(-4)$-spheres to $b^+$.
\item \emph{Why the manuscript reads as alien.}
\begin{itemize}
\item Theorem~\ref{thm:FL} and Lemma~\ref{lem:adj} are never stated as such.
\item Classical objects are replaced by bespoke ones: holonomy sweeps for $\mu$-classes, explicit toric estimates for wall-crossing, case-by-case dimension counts for one inequality.
\item The argument is run on the capped 4-manifold rather than in Floer homology.
\end{itemize}
\item \emph{A Floer-theoretic (DLME-shaped) proof.} One would want an $S^1$-equivariant $\mathrm{SO}(3)$-monopole Floer theory of $(Y_{\pm2},t)$, spin-graded, with localization weights $2^{n_D-1}$, in which $w$ lowers $\ell_t$ by $\tfrac18$ per period.
\begin{itemize}
\item Monotonicity cannot hold for single cobordisms: $HB\simeq1$ mod 2 has spin shift $-\tfrac12$.
\item For blocks with irreducible Floer ends it fails because of end terms at the 3-dimensional Seiberg--Witten points. $Y_2$ is not an L-space, since $2<2g-1$. Only a version including the central flats is viable, at 20--35\%.
\item No such theory exists. The capped construction is precisely the device that avoids it (\texttt{notes/B2}, \texttt{notes/B3}).
\end{itemize}
\item \emph{Consistency checks.}
\begin{itemize}
\item At $\pm1$ the same scheme reduces to DLME's energy argument, since $\kappa<0$ \V.
\item For $\pm p$ with $p\ge3$ it fails at the outset, because of non-central $U(1)$ flats \V.
\item With spacing $1$ (no cosmetic map) the exclusions provably fail, as they must.
\end{itemize}
\item \emph{Risk.} Overall confidence that the reorganized argument is a correct proof is 35--45\%. The weak points are:
\begin{itemize}
\item the projection for mixed limits;
\item uniformity of the chamber condition at the faces of the cube;
\item persistence of the $\mu(S)$ representatives under Uhlenbeck limits, which Donaldson's divisor $V_S$ should settle;
\item the lens cancellation in the coupled problem.
\end{itemize}
\end{itemize}
\end{document}
